\section{Proofs}
\label{app:proofs}

\subsection{Projective and diamond distance}

\begin{lemma}[Unitary-channel diamond comparison]
\label{lem:diamond-comparison}
For the unitary channels
$\mathcal U(X)=UXU^\dagger$ and $\mathcal V(X)=VXV^\dagger$, define the
\emph{unnormalized} diamond distance
$d_\diamond(\mathcal U,\mathcal V)
:=\|\mathcal U-\mathcal V\|_\diamond\in[0,2]$.  Then
\begin{equation*}
 d_{\mathrm{proj}}(U,V)
 \leq d_\diamond(\mathcal U,\mathcal V)
 \leq 2d_{\mathrm{proj}}(U,V).
\end{equation*}
If the normalized convention
$\widetilde d_\diamond:=\tfrac12\|\mathcal U-\mathcal V\|_\diamond$ is used,
the corresponding explicit relation is
$\tfrac12d_{\mathrm{proj}}\leq\widetilde d_\diamond\leq
d_{\mathrm{proj}}$.
\end{lemma}

\begin{proof}
Let $\Theta\in[0,2\pi]$ be the length of a shortest closed arc containing
the spectrum of $U^\dagger V$.  Centering that arc by a global phase puts
every eigenvalue within angle $\Theta/2$ of $1$; conversely, for every
phase the rotated spectrum still spans an arc of length $\Theta$, so some
eigenvalue lies at angle at least $\Theta/2$ from $1$.  Hence
\begin{equation*}
 d_{\mathrm{proj}}(U,V)=2\sin(\Theta/4).
\end{equation*}
The numerical-range formula for two unitary
channels~\cite{watrous2018theory} gives
\begin{equation*}
 d_\diamond(\mathcal U,\mathcal V)=
 \begin{cases}
  2\sin(\Theta/2),&0\leq\Theta\leq\pi,\\
  2,&\pi\leq\Theta\leq2\pi.
 \end{cases}
\end{equation*}
For $\Theta\leq\pi$, the ratio is $2\cos(\Theta/4)\in[\sqrt2,2]$;
for $\Theta\geq\pi$, one has
$d_{\mathrm{proj}}\in[\sqrt2,2]$ and $d_\diamond=2$.  The two inequalities
follow.
\end{proof}

\subsection{The accuracy grid}

\begin{proof}[Proof of \cref{cor:w1-epsilon-cyclic}]
At $\epsilon=1/8$ one has $Q_\epsilon=6$, so $Q_\epsilon\geq6$ on the
whole range.  The bounds $2\epsilon\leq\arcsin(2\epsilon)\leq\pi\epsilon$
and $\lfloor a\rfloor\geq a-1$ give
$1/(4\epsilon)\leq1/(2\epsilon)-1\leq K_\epsilon\leq\pi/(4\epsilon)$.
Since $Q_\epsilon\leq\pi/[2\arcsin(2\epsilon)]$, we have
$\pi/(2Q_\epsilon)\geq\arcsin(2\epsilon)$, and
\cref{thm:w1-exact-separation} with $K=Q_\epsilon$ gives the separation.
For the count, fix a non-identity Pauli $P$.  Both characters $\pm1$ occur,
so, up to a global sign, $R_P(\theta)R_P(\theta')^\dagger$ has eigenvalue
arguments $\mp g/2$, where $g\in[0,\pi]$ is the circular distance between
$\theta$ and $\theta'$ in $\mathbb R/2\pi\mathbb Z$, and centering gives $d_{\mathrm{proj}}=2\sin(g/4)$.  With
$g_\epsilon:=4\arcsin(2\epsilon)$ one has $2\sin(g_\epsilon/4)=4\epsilon$,
so pairwise $4\epsilon$-separation is pairwise circular distance at least
$g_\epsilon$.  Points on a circle of circumference $2\pi$ with pairwise
gaps at least $g_\epsilon$ number at most $\lfloor2\pi/g_\epsilon\rfloor$,
and equal spacing attains this; finally
$\lfloor2\pi/g_\epsilon\rfloor=\lfloor\pi/(2\arcsin(2\epsilon))\rfloor
=Q_\epsilon$.
\end{proof}

\subsection{Output-code modes}

The codes of \cref{def:info-budget-bits} are fixed prefix-free streaming
binary codes.  $E_{\mathrm{self}}$ serializes a target-basis circuit with
gate opcodes, all qubit and generator addresses, record delimiters, and
sparse parameters, and assumes no dictionary.  Given
$\mathcal F_m=(m,P_1,\ldots,P_m,K,\Delta)$,
$E_{\mathrm{rel}}(\cdot\mid\mathcal F_m)$ has four modes.  (i) Aggregate
mode: a self-delimiting header declaring a block length $b$, followed by
the $m$ aggregate digits packed in base $K$ in consecutive blocks of $b$
digits, each written least significant bit first as a
$\lceil b\log_2K\rceil$-bit integer, and a final block of $b'<b$ digits in
$\lceil b'\log_2K\rceil$ bits.  (ii) Literal mode: complete target-basis
records.  (iii) Ordered-share mode: a header declaring a split point,
after which the fixed public schedule lets generator addresses be omitted
from the residue slots, packed in base $Q$, while round-boundary and
packed-block delimiters are retained.  (iv) $p$-block hybrid mode: a
header declaring the pass blocks and passthrough coordinates, followed by
packed aggregate blocks and, for each round, the passthrough residues of
that round packed in base $Q$, in the order in which one left-to-right
scan per pass produces them, a public function of $(m,p,q,r)$.


\subsection{Upper bounds}

\begin{proof}[Proof of \cref{prop:w1-block-hybrid}]
Partition the first $q$ coordinates into at most $p$ consecutive blocks of
size at most $h$.  Each pass is re-delivered by the input source
(\cref{def:w1-pass-accounting}), so the compiler keeps no copy of the
stream between passes.  On pass $t$, keep only the residues of block $t$ and add
their first $r-1$ rounds modulo $Q$; when the final-round mates arrive,
emit the corresponding aggregate base-$Q$ block.  The partial sums are held
as one packed integer in the store field at its literal bit
length (\cref{sec:w1-dispersed}); the live packed block, indices, and framing
give \eqref{eq:w1-upper-state} at the crossings.  Each block is packed
separately, which
costs at most one bit per block; together with the per-block framing this is
the $O(p\log(mQrp))$ term in the output bound.  On the first pass,
emit every share of each unselected coordinate in the declared round-major
order, packing the residues of each round in base $Q$ and flushing them
at the end of that round, so that the residues of round $r-1$ are written
before the crossing; these rotations already have the correct product.  The selected coordinates can be balanced
across the $p$ blocks, so the maximum block has size $\lceil q/p\rceil$.
Their consecutive boundaries are a public function of $(m,p,q)$, so the
header stores those integers rather than $p$ separate delimiters.
The prefix contains $r-1$ residues and the
whole output contains $r$ residues for each such coordinate, giving
\eqref{eq:w1-upper-A} and \eqref{eq:w1-upper-output}.  If $ph\geq m$ there
are no passthrough coordinates, $A_p$ is the header, and
$p\lceil m/p\rceil\log_2Q\leq m\log_2Q+p\log_2Q$ gives the stated bound on
$A_p+pS$ for any $r$.  For $r=2$ and general $q$, $(m-q)\log_2Q+p\lceil
q/p\rceil\log_2Q$ is at most $m\log_2Q+p\log_2Q$, which gives the same
bound; since every forward snapshot is bounded by the same
block, the bound also holds for $A_p+\Sigma^{\to}_p$.
\end{proof}

\begin{proof}[Proof of \cref{prop:w1-semantic-upper}]
Buffer the incoming aggregates $b$ at a time, held as one packed integer,
and write each full buffer as one $\lceil b\log_2K\rceil$-bit integer and
a final partial block of $b'<b$ digits in $\lceil b'\log_2K\rceil$ bits;
the buffer, an index, and framing are the only live state.  Each block loses less than one bit to rounding,
which gives the $\lceil m/b\rceil$ term.  For the single packed integer,
let $P=\sum_{i\leq j}x_iK^{i-1}$ be the value of the first $j$ digits and
$N=P+K^jR$ the final integer.  Since $K^jR\equiv0\pmod{2^{ej}}$ and the
higher bits of $N$ depend on $R$, a correct write-once writer has written at
most the low $ej$ bits of $P$, so at most $2^{ej+1}$ distinct written
prefixes occur.  Two digit prefixes with the same snapshot and the same
written prefix produce the same $N$ for every $R$, hence $P=P'$.  The
$K^j$ values of $P$ therefore need at least $K^j/2^{ej+1}=K_{\mathrm{odd}}^{\,j}/2$
distinct snapshots, which, being self-delimiting
(\cref{sec:w1-dispersed}), need at least
$j\log_2K_{\mathrm{odd}}-1$ bits by Kraft's inequality; take
$j=\lfloor m/2\rfloor$.  \Cref{prop:w1-stream-equivalence} gives
semantic equality with the flat representation.  Prefix the output with the
abort-coin branch tag for the randomized construction and take
expectations.  The self-contained construction writes each nonzero
$R_{P_j}(2\pi x_j/Q)$ with its address and rational parameter.
\end{proof}

\subsection{Physical streams}

\begin{proof}[Proof of \cref{lem:clifford-interleaving}]
Write the emitted unitary as $D_r'L_r\cdots D_1'L_1$ with $L_t$ the
round-$t$ layer and $D_t'$ the Clifford record following it, so
$D_t=D_t'\cdots D_1'$.  Moving each Clifford leftwards replaces $L_t$ by
$D_{t-1}^\dagger L_tD_{t-1}$ and leaves $D_r$ as a suffix.  Since
$L_t=\prod_jR_{P_j}(\Delta s_j^{(t)}a_j^{(t)})$ and conjugation is an
automorphism, the hypothesis gives
$D_{t-1}^\dagger L_tD_{t-1}
=\prod_jR_{P_{\pi_t(j)}}(\Delta\varepsilon_j^{(t)}s_j^{(t)}a_j^{(t)})$,
which is a round of a scheduled stream after the stated
re-indexing.  The Clifford records lie in the stream, so $D_{t-1}$, and
hence $\pi_t$ and $\varepsilon^{(t)}$, are computable from the prefix.
Finally $D_r$ is common to every member of the family and unitary, so it
preserves projective distance and the separation of
\cref{thm:w1-exact-separation} is unaffected.
\end{proof}

\begin{proof}[Proof of \cref{prop:physical-instantiation}]
In (a) and (b) the stated coefficient-freedom hypothesis makes the round
coefficients unconstrained across rounds,
so every residue sequence---in particular every valid share
decomposition of \cref{def:w1-dispersed-stream}---arises from some
schedule, and \cref{prop:w1-stream-equivalence} supplies realizability
and token privacy verbatim.  For prefixes: after any prefix that ends
before the last round, a final-round residue choice reaches every
aggregate in $\mathbb Z_Q^m$.  In (c), sign normalization $a\mapsto sa$ is
a bijection of $\mathbb Z_Q$ for either sign, so normalizing every
record maps the signed family onto the unsigned family
record-for-record, and the same properties transfer; the distributional
clause is the one proved in \cref{prop:w1-stream-equivalence}, applied
after normalization.  For the constant schedule, $x_j=ra_j\bmod Q$ is a
function of the first round.  An exhaustive enumeration at
$(m,r,Q)=(2,2,5)$, frozen at
\path{data/frozen/fooling_set_m2_r2_Q5.json}, independently
confirms all three behaviors: every time-varying and every signed
prefix class reaches all $Q^m$ aggregates, and each constant-schedule
prefix class reaches exactly one.
\end{proof}

\begin{proof}[Proof of \cref{prop:frame-cross-cut}]
(i) Sign normalization is a bijection of $\mathbb Z_Q$ for either sign, so
normalizing every record maps the twirled family onto the family of
\cref{def:w1-dispersed-stream} while preserving aggregates, record
boundaries, and cut positions; each normalized stream is the image of the
twirled streams that differ only in their sign patterns.  The cited
theorems apply to the image and pull back verbatim.

(ii) Take prefixes whose residue fields emit the residue $1$ exactly
once on each generator and are otherwise zero, and whose sign fields
realize the frame class $\sigma\in\{\pm1\}^m$; the residue fields are
identical across classes.  Choose one common suffix adding a residue
$a_j$ on generator $j$, for instance $a_j=0$.  The completed aggregates
$x_j=\sigma_j+a_j\bmod Q$ then differ between two classes in every
coordinate where the frames differ, because $2\not\equiv0\pmod Q$ for
$Q\geq3$, and all completions are valid family instances.  By
\cref{thm:w1-exact-separation} the corresponding targets are separated by
$2\sin(\pi/(2Q))>2\epsilon$.  For a deterministic one-pass compiler the
suffix computation is a function of the snapshot and the common suffix;
a run that halted before $c$ would give the same output for two classes
with different targets, so every run reaches $c$, and the pair (prefix
output, snapshot) must be injective on the $2^m$ frame classes; prefix-freeness of the self-delimiting snapshot followed by
the committed prefix and Kraft's inequality give the deterministic bound.
For the multipass statement, draw $\sigma$ uniformly and repeat the
forward-transcript argument of \cref{thm:w1-unified-frontier} with the
common suffix in place of $V$: the forward transcript determines the
final output, the decoder recovers $\sigma$ with error at most $\delta$,
and Fano's inequality gives $I(\sigma;T)\geq(1-\delta)m-h_2(\delta)$.
The frame register is classical restart state, so its bits lie in the
named fields of the snapshot.

In folded mode the sign is consumed at emission, so the crossing
information consists of the signed residues $\pm1$ themselves, which are
distinct for $Q\geq3$, and the argument applies verbatim to the folded
records.
\end{proof}

\begin{proof}[Proof of \cref{prop:eta-pareto}]
Condition on the mask $T$.  Restrict the signs to $s\equiv+1$ and fix the
exposed coordinates; $T$ and the exposed values are then public, known to
both Alice and Bob in the transcript argument, and are built into the
transition function rather than stored as state.  The targets are the
dispersed family on the $|T|$ generators in $T$ multiplied by one fixed
unitary $W$ of the exposed coordinates.  The proof of
\cref{thm:w1-unified-frontier} uses only pairwise separation of the
targets, which this translation preserves, so it applies unchanged and
gives
$\overline A_p(T)+pS(T)\geq(1-\delta)|T|\log_2K-h_2(\delta)
-\lceil\log_2(p+1)\rceil$ for every fixed $T$.  Take the expectation over
the Bernoulli mask and use $\mathbb E|T|=\eta m$.
The one-pass clause follows in the same way from
\cref{thm:w1-one-pass-mask} applied to the dispersed coordinates.
\end{proof}

\subsection{Commitment toll}

\begin{proof}[Proof of \cref{cor:unitary-toll}]
For $j\in T$ the signed residue $c_j:=s^{(1)}_ja^{(1)}_j\bmod Q$ is
uniform on $\mathbb Z_Q$, because $a^{(1)}_j$ is uniform and independent
of $s^{(1)}_j$; so $(T,c)$ is uniform on $\binom{[m]}{k}\times\mathbb Z_Q^k$.
If $y\neq y'$, the corresponding round-one actions differ by a nonzero
multiple of $\Delta=2\pi/Q$ on some generator, and
\cref{thm:w1-exact-separation} with the full residue alphabet $K=Q$
separates them by $2\sin(\pi/(2Q))>2\epsilon$; integer and modular
residues give the same rotations up to a global phase $\pm1$.  No prefix is within $\epsilon$ of both, so
an executed prefix determines $y$, and mask uniformity makes the
prefixes one prefix-free code; source coding gives
$\mathbb E[A_{\mathrm{pre}}]\geq H(y)$.  The map $(T,c)\mapsto(y,T)$ is
injective, so $H(y)=H(T,c)-H(T\mid y)$.  Given $y$ with support of
size $\nu$, the consistent masks are the $k$-sets containing the
support, all equally likely, so
$H(T\mid y)=\mathbb E\log_2\binom{m-\nu}{k-\nu}$ with
$\nu=|\operatorname{supp}y|\sim\mathrm{Bin}(k,1-1/Q)$; this is
\eqref{eq:unitary-toll}.  Equivalently, writing $V:=\operatorname{supp}y$,
the values of $y$ on $V$ are uniform on $(\mathbb Z_Q\setminus\{0\})^V$ and
independent of $T$ given $V$, so $H(T\mid y)=H(T\mid V)$ and the excess of
$H(y)$ over $k\log_2Q$ is $H(T)-H(T\mid V)=I(T;V)$.  The deferred compiler stores the $k$ signed
residues in the public round-major order and recovers $T$ from the
round-two record types, as in \cref{prop:commitment-toll}(i); the target
depends on the signs only through the signed aggregate
\eqref{eq:scheduled-aggregate}, so the folded residues suffice.
Subtracting the deferred state gives the toll.  For
\eqref{eq:unitary-toll-rate}, use $I(T;V)=H(V)-H(V\mid T)$.  Given $T$,
each member of $T$ lies in $V$ independently with probability $1-1/Q$,
so $H(V\mid T)=k\,h_2(1/Q)$.  Each $j\in[m]$ lies in $V$ with
probability $\eta(1-1/Q)$, so subadditivity gives
$H(V)\leq m\,h_2(\eta(1-1/Q))$.  Conversely $V$ is a uniformly random
subset of $[m]$ given its size $\nu\sim\mathrm{Bin}(k,1-1/Q)$, so
$H(V)\geq\mathbb E\log_2\binom{m}{\nu}$.  The bound
$\log_2\binom ms\geq m\,h_2(s/m)-\log_2(m+1)$, a second-order expansion
of $h_2$ around $\eta(1-1/Q)$, which lies in $(0,1)$, and a Chernoff tail
bound for $\nu$ give $H(V)\geq m\,h_2(\eta(1-1/Q))-O(\log m)$.  Dividing by
$k=\eta m$ gives the limit, approached from below.  For attainability,
given $\nu$ the support is uniform over the $\nu$-subsets of $[m]$ and the
nonzero values are uniform on $(\mathbb Z_Q\setminus\{0\})^\nu$, so the
stated prefix has expected length
$H(y)-H(\nu)+O(\log m)\leq H(y)+O(\log m)$, and it implements the round-one
action on its own.  Strict concavity of
$h_2$ with $h_2(0)=0$ gives $h_2(\eta u)>\eta h_2(u)$ for $0<\eta<1$ and
$u=1-1/Q$, which is positivity, since $h_2(u)=h_2(1/Q)$.  The derivative
in $u$ is $h_2'(\eta u)-h_2'(u)>0$ because $h_2'$ is decreasing, so the
rate increases with $Q$, to $h_2(\eta)/\eta$ as $u\to1$.
\end{proof}
